\section{Reduced-product composition and bounded-power laws}
\label{sec:p03-fv-reduced-annex}

\begin{annexitem}{annex:p03-fv-finite-partition-lift}{thm:p03-fv-finite-partition-lift}{An exact lifted partition}
\begin{proof}
If \(n=1\), take \(b_1=1_B\). Otherwise choose, for the finite
list \(k<n\), preimages \(u_k\) of \(c_k\). Set
\[
 b_1=u_1,\qquad
 b_k=u_k\wedge\neg\bigvee_{j<k}b_j\ (1<k<n),\qquad
 b_n=\neg\bigvee_{j<n}b_j.
\]
Every two cells are disjoint and their join is one. Induction gives
\(h(b_k)=c_k\) for \(k<n\), since \(c_k\) is disjoint from the
earlier cells; the last image is their complementary cell \(c_n\).
This also covers zero cells and trivial quotient algebras. Finite choice
requires no independent axiom beyond classical ZF. Apply the construction
to the surjective quotient map from \(\mathcal P(\Omega)\), where a
Boolean partition is literally a partition into subsets of \(\Omega\).
\end{proof}
\end{annexitem}

\begin{annexitem}{annex:p03-fv-reduced-well-defined}{thm:p03-fv-reduced-well-defined}{Representative invariance and term evaluation}
\begin{proof}
Reflexivity and symmetry of element equivalence are immediate. If \(a,b\)
disagree only on \(E\in I\) and \(b,c\) only on \(D\in I\),
then \(a,c\) disagree only on a subset of \(E\cup D\), which belongs
to \(I\). For a finite tuple of equivalent input representatives their
values agree outside the union of their finitely many disagreement sets.
Each coordinate function has equal output on that complement; each relation
has identical truth there. Function outputs are therefore equivalent, and
the two relation truth sets have symmetric difference in \(I\). Membership
of their complements in \(I\) is the same, establishing relation invariance.
Constants are their declared coordinate tuples. Choice and nonempty factor
sorts make every product sort nonempty, hence every quotient sort nonempty.

For an arbitrary factor formula \(\psi\), outside the same union the
parameter tuples are identical in the very same factor. Their satisfaction
of \(\psi\) is thus identical, regardless of its bound quantifiers.
Consequently its two local truth sets give the same Boolean class. This
uses equality of factor parameters, not an invalid formula transfer theorem.
Term induction starts at variables and coordinate constants; the function
case follows by the just established descended interpretation. Hence term
values are coordinatewise classes. Equality of two term classes is exactly
that their disagreement set lies in \(I\), or that their equality truth
set has class one. Relation atoms have the same class-one criterion by
definition, including nullary atoms. This proves every asserted clause.
\end{proof}
\end{annexitem}

\begin{annexitem}{annex:p03-fv-reduced-reduction}{thm:p03-fv-reduced-reduction}{The complete quotient-Boolean induction}
\begin{proof}
Use induction on the finite syntax of the input formula. Atomic satisfaction
is the class-one criterion just proved. Logical falsity is false in the
quotient structure and \(0\ne0\) is false in its Boolean algebra.
For negation, negate the induction equivalence as a statement of ordinary
two-valued satisfaction; this negates \(\Theta\), not one of its argument
classes. For conjunction use the two inductive equivalences, concatenated
factor lists and the specified disjoint renaming. These are precisely the
given reduction clauses. The expansions of other connectives and universal
quantification consequently preserve their semantics.

Consider \(\exists y_s\chi(\bar x,y_s)\). Suppose first that a class
\(\pi_I(b)\) witnesses its matrix. At every index the finite list of
factor-formula truths at \((\bar a(\gamma),b(\gamma))\) determines a
unique pattern \(\sigma\). Let \(S_\sigma\) be the set of indices
with that pattern. They partition \(\Omega\) exactly. Their classes
are a Boolean partition; each \(S_\sigma\) is a subset of the availability
set \(X_{\eta_\sigma}(\bar a)\), and
\[
 X_{\psi_j}(\bar a,b)=\bigcup_{\sigma_j=1}S_\sigma.
\]
The matrix induction hypothesis therefore supplies the final \(\Theta\)
condition and proves the existential Boolean clause.

Conversely assume that clause holds, with Boolean cells \(Z_\sigma\).
Use finite partition lifting to choose actual sets \(S_\sigma\)
partitioning \(\Omega\), with \([S_\sigma]_I=Z_\sigma\).
The Boolean inequality to availability means
\[
 E_\sigma:=S_\sigma\setminus X_{\eta_\sigma}(\bar a)\in I.
\]
There are finitely many patterns, so \(E=\bigcup_\sigma E_\sigma\)
belongs to \(I\). At every \(\gamma\notin E\), the unique cell
\(S_{\sigma(\gamma)}\) supplies a nonempty set of actual witnesses
\[
 \{d\in(A_\gamma)_s:
       A_\gamma\models\delta_{\sigma(\gamma)}(\bar a(\gamma),d)\}.
\]
At every \(\gamma\in E\), use the nonempty whole sort as the available
set. Choice selects from all these sets and gives an actual full product
element \(b\). Outside \(E\) it realizes exactly the chosen complete
pattern, including every negative bit. It follows that
\[
 X_{\psi_j}(\bar a,b)\mathbin{\triangle}
                \bigcup_{\sigma_j=1}S_\sigma\subseteq E.
\]
Thus its local class is exactly \(\bigvee_{\sigma_j=1}Z_\sigma\).
The assumed last Boolean condition and matrix induction now give
\(A_I\models\chi(\pi_I(\bar a),\pi_I(b))\). This is the required
quotient witness. It is not necessary to realize the prescribed patterns
on \(E\), nor legitimate to discard \(E\) without first proving it
belongs to the ideal.

When \(m=0\), the one pattern is the empty conjunction. Its availability
set is all indices by nonemptiness of the quantified sorts. The same proof
selects the required witness, while the zero-input Boolean condition stays
unchanged. Every syntactic recursion step produces a finite list and a
finite Boolean formula from proper subformulas; the partition has \(2^m\)
cells. This establishes effective termination. When \(I=\{\varnothing\}\)
every exceptional set above is empty and the quotient algebra is the full
index algebra, recovering the direct-product proof.
\end{proof}
\end{annexitem}

\begin{annexitem}{annex:p03-fv-reduced-elementarity}{thm:p03-fv-reduced-elementarity}{Parameter transfer through the explicit condition}
\begin{proof}
Coordinate inclusion takes equivalent representatives to equivalent ones,
and reflects equivalence because it preserves equality of their values.
It hence induces an injective structure map on quotient sorts. For a finite
parameter tuple from the smaller product and any factor formula \(\psi\),
factorwise elementarity makes the two local truth sets identical. The
composition theorem evaluates every product-language formula by the same
Boolean formula at exactly the same classes in the same Boolean algebra.
All such formulas agree, proving elementary inclusion. Given specified
elementary embeddings instead, coordinate equality is reflected and each
local formula is preserved and reflected, so the same argument applies.

For elementary-equivalent, rather than included, factors, take an input
sentence. The construction's component formulas have no free variables.
Their truth in \(A_\gamma\) and \(B_\gamma\) therefore agrees at
every index. Their truth sets are identical, and evaluating the computed
Boolean sentence proves equality of the two reduced-product theories.
No representative map between their carriers is needed or asserted.
\end{proof}
\end{annexitem}

\begin{annexitem}{annex:p03-fv-reduced-example}{thm:p03-fv-reduced-example}{The two witness cells and an explicit separating filter}
\begin{proof}
For the matrix choose the local list \(R(x,y),P(y)\). Its Boolean
condition is \(X_R=1\land X_P\ne1\). In the existential partition,
the cells with an \(R\)-false bit must be zero. The two remaining cells
\(Z_{10},Z_{11}\) must be disjoint, join to one, lie under \(U,V\)
respectively, and satisfy \(Z_{10}\ne0\). These conditions imply
\(U\vee V=1\) and \(U\ne0\). Conversely those two conditions allow
\(Z_{10}=U\) and \(Z_{11}=\neg U\), since \(\neg U\le V\);
the other cells are zero. The composition theorem proves that this Boolean
equivalence is the exact satisfaction criterion.

In the displayed cofinite-filter example, every factor tuple has \(P\)
truth set exactly the odd numbers, independent of its chosen elements.
Its complement, the even numbers, is infinite, so the relation \(P\)
is false in the reduced product. Its logical negation is therefore true.
But the even numbers are not cofinite; the proposed unrestricted \L o\'s
criterion for that negation would say it is false. For the worked existential,
the local \(R\) relation is universal; availability for the negative
\(P\) pattern is the even numbers and for the positive pattern is the
odd numbers. Their classes join to one, and the even class is nonzero,
so the criterion holds, as the actual witness semantics also shows.
For an ultrafilter every set class is zero or one and the quotient Boolean
algebra is nontrivial; the two classes in this example could not both be
nonzero there. This identifies the exact missing ultrafilter property.
\end{proof}
\end{annexitem}

\begin{annexitem}{annex:p03-fv-source-power-products}{thm:p03-fv-source-power-products}{Constants and both product isomorphisms}
\begin{proof}
For a nontrivial \(B\), BPIT supplies its ultrafilters, so \(B^*\)
is nonempty. Every constant function is continuous. Evaluation at any point
therefore hits every base value, establishing the subdirect property, and
distinct constants are distinct functions, establishing the base embedding.
The two-element Boolean algebra has just its unique proper ultrafilter;
evaluation at that one point is an isomorphism \(A[\mathbf2]^*\cong A\).
The trivial algebra has no proper ultrafilters, so there is just the empty
function in its power; this proves the stated exceptional case.

For the coefficient-product assertion, an ultrafilter of \(B_1\times B_2\)
contains exactly one of \((1,0),(0,1)\). In the first case it has the
form \(\{(b_1,b_2):b_1\in U_1\}\) for a unique ultrafilter \(U_1\)
of \(B_1\); in the second case the analogous \(U_2\) description
holds. To verify the first description, membership of \((1,0)\) makes
membership of \((b_1,b_2)\) equivalent to that of \((b_1,0)\),
and the ultrafilter axioms restrict to those of \(U_1\); the converse
construction has the same axioms. The second case is symmetric. Basic
clopens restrict to the corresponding basic clopens in the two spaces.
Thus the Stone space is their disjoint union, each component clopen.
Restriction gives a bijection
\[
 C(B_1^*\sqcup B_2^*,A)\longrightarrow
                       C(B_1^*,A)\times C(B_2^*,A):
              c\longmapsto(c|_{B_1^*},c|_{B_2^*}).
\]
Its inverse combines the two functions on their disjoint components;
clopenness proves continuity in both directions. Restriction commutes
pointwise with every finitary operation, including constants, so it is an
algebra isomorphism. Empty components cause no change to this argument.

For the base-product assertion, a continuous function into the discrete
\(A_1\times A_2\) has continuous coordinate functions. Conversely two
continuous coordinate functions give a continuous pair function: each
fiber at \((a_1,a_2)\) is the intersection of the two coordinate fibers,
and the preimage of an arbitrary subset is their union. Alternatively their
finite clopen partitions have a common finite refinement. These operations
are inverse bijections. Each basic operation of the algebra product acts
coordinatewise, so the bijection commutes with it at every Stone point and
is an isomorphism. This proves all of the source's product assertions with
their stated domain qualifications.
\end{proof}
\end{annexitem}

\begin{annexitem}{annex:p03-fv-iterated-power}{thm:p03-fv-iterated-power}{Why currying preserves finite range and continuity}
\begin{proof}
Put \(X=B_1^*\), \(Y=B_2^*\). A continuous map from compact
\(Y\) to the discrete algebra \(C(X,A)\) has finite range
\(f_1,\ldots,f_n\) and clopen fibers \(V_1,\ldots,V_n\).
The uncurried map equals \(f_k(x)\) on \(X\times V_k\).
For each \(a\in A\) its fiber is
\(\bigcup_{k=1}^n f_k^{-1}(a)\times V_k\), which is clopen.
The finitely many finite ranges of the \(f_k\) have finite union.
It is hence continuous into discrete \(A\).

Conversely let \(f:X\times Y\to A\) be continuous. If either space
is empty the uncurried domain is empty and both constructions have exactly
one element; this handles those cases separately. Otherwise product
compactness and discreteness give finitely many clopen fibers of \(f\).
For each point, one such fiber is an open neighborhood. Since products of
clopen sets are a basis for the product of Stone spaces, that point lies
in a rectangle \(U\times V\) on which \(f\) is constant. Compactness
of \(X\times Y\) gives a finite covering list of such rectangles
\(U_k\times V_k\). Take the finite Boolean partition of \(Y\)
generated by the finitely many \(V_k\). For \(y,y'\) in one nonempty
cell, membership in each \(V_k\) agrees. For any \(x\), a covering
rectangle containing \((x,y)\) therefore also contains \((x,y')\),
so \(f(x,y)=f(x,y')\). Its continuous slices \(x\mapsto f(x,y)\)
are thus constant as elements of the discrete function algebra on each
of finitely many clopen \(Y\)-cells. The curried map has finite range
and clopen fibers, so is continuous with the required discrete codomain.
Currying and uncurrying are inverse and commute pointwise with operations.

Stone representation identifies the compact Stone space of
\((X\times Y)^*\) with \(X\times Y\), making the first claimed
isomorphism. Interchange of the two product coordinates is a homeomorphism
and commutes with pointwise operations, giving the symmetric iterated-power
isomorphism. Finite rectangle covers, not arbitrary infinite patching, are
what justify the boundedness assertion.
\end{proof}
\end{annexitem}

\begin{annexitem}{annex:p03-fv-power-quotient}{thm:p03-fv-power-quotient}{Coefficient transport, its kernel, and actual lifted preimages}
\begin{proof}
First, for \(b\in B\), the basic clopen \(N_b\) contains \(F^*\)
if \(b\in F\). If \(b\notin F\), the filter generated by
\(F\cup\{\neg b\}\) is proper: a zero intersection would put some
\(f\in F\) below \(b\), forcing \(b\in F\). BPIT extends this
proper filter to an ultrafilter containing \(F\) and \(\neg b\).
It lies outside \(N_b\). Thus the equalizer containment formulation of
\(\theta_F\) is exactly membership of its Boolean coefficient in \(F\).
For completeness, put \(N_F=\{b\in B:\neg b\in F\}\). It is an
ideal: zero belongs, passing to a smaller element enlarges its complement,
and the complement of a binary join is the intersection of the complements.
Agreement modulo \(F\) is equivalently that the Boolean symmetric
difference \(\Delta(b,c)=(b\wedge\neg c)\vee(\neg b\wedge c)\)
belongs to \(N_F\). To prove transitivity, split
\(b\wedge\neg d\) by \(c\vee\neg c=1\): its \(c\)-part
lies below \(c\wedge\neg d\), and its \(\neg c\)-part below
\(b\wedge\neg c\). Splitting \(d\wedge\neg b\) in the same
way gives \(\Delta(b,d)\le\Delta(b,c)\vee\Delta(c,d)\).
Closure of the ideal proves transitivity; reflexivity and symmetry follow
from \(\Delta(b,b)=0\) and \(\Delta(b,c)=\Delta(c,b)\).
For operations, distributivity gives
\[
 (b\wedge d)\wedge\neg(c\wedge e)
   =(b\wedge d\wedge\neg c)\vee(b\wedge d\wedge\neg e)
   \le(b\wedge\neg c)\vee(d\wedge\neg e).
\]
The reversed difference has the analogous bound with the pairs reversed,
so \(\Delta(b\wedge d,c\wedge e)\le
\Delta(b,c)\vee\Delta(d,e)\). Also
\(\Delta(\neg b,\neg c)=\Delta(b,c)\). De Morgan's identity and
the meet inequality therefore give the same bound for the join operation.
These bounds prove congruence for meets, joins and complements; constants
have no representative choice. Hence all Boolean operations descend.
Thus \(B/F\) is a Boolean algebra and \(q_F\) is a surjective Boolean
homomorphism. For a proper \(F\) its zero and one classes differ; for
\(F=B\) they agree and the quotient is the trivial Boolean algebra.
The quotient map has
\(q_F(b)=1\Longleftrightarrow b\in F\) and
\(q_F(b)=0\Longleftrightarrow\neg b\in F\).

Use the formal finite-partition presentation. Coefficient transport along
\(q_F\) preserves joins, meets, zero and one, so sends each presentation
to a legitimate partition and sends equivalent presentations to equivalent
ones. It respects all operation refinements and constants, defining an
algebra homomorphism to \(A[B/F]_{\mathrm{part}}\).
For two presentations, refine to cells carrying paired labels. They have
equal images exactly when the join of cells with unequal labels has image
zero. Its complementary equal-label join has image one, equivalently is
in \(F\). That is exactly \(\theta_F\). The kernel of a homomorphism
is a congruence: reflexivity, symmetry and transitivity follow from equality
of images, and preservation of operations proves compatibility.

Given any presentation over \(B/F\), ensure it has at least one cell.
If it is empty, then \(B/F\) is trivial; replace it by its equivalent
one-cell presentation with coefficient \(0=1\) and any base label.
Lift its finite coefficient
partition by Theorem~\ref{thm:p03-fv-finite-partition-lift}, leaving
each label unchanged. This supplies an actual preimage presentation over
\(B\), proving surjectivity. Passing to kernel classes gives a bijective
operation-preserving map, whose inverse preserves operations because every
tuple has preimages. It is the asserted algebra isomorphism. Finally
transport through the proved partition/continuous-power identifications
on both sides to obtain the source's starred-power formulation. For the
improper filter every equalizer contains the empty set \(F^*\), so the
congruence is universal and the codomain power over the trivial quotient
has one element. This also verifies the claimed improper-filter boundary.
\end{proof}
\end{annexitem}

\begin{annexitem}{annex:p03-fv-power-cardinality}{thm:p03-fv-power-cardinality}{Both cardinal bounds}
\begin{proof}
Let \(\kappa=\max(|A|,|B|)\), which is infinite under the assumptions.
Every power element has a finite labeled-partition presentation, and the
valid presentations form a subset of the finite lists from \(B\times A\).
Their quotient maps onto the power carrier.
Under Choice that list set has cardinal at most
\(\sum_{n<\omega}(|B|\cdot|A|)^n=\kappa\).
This proves the upper bound without asserting that every list is a valid
presentation.

Constants inject \(A\) into the power, since \(B\) is nontrivial.
Fix distinct \(a_0,a_1\in A\). To \(b\in B\) associate the
partition with coefficient \(b\) labeled \(a_1\) and coefficient
\(\neg b\) labeled \(a_0\). Equality of two such presentations
forces both \(b\wedge\neg c=0\) and \(\neg b\wedge c=0\),
which is exactly \(b=c\). Thus \(B\) also injects into the power.
The two lower bounds give at least \(\kappa\); together with the
upper bound they prove equality. Infinite cardinal multiplication under
Choice gives \(|A|\cdot|B|=\kappa\) as stated. For a singleton
base all labels are equal and all presentations are equivalent, proving
the asserted failure when its nontriviality hypothesis is deleted.
\end{proof}
\end{annexitem}
